\documentclass[11pt]{article}
\usepackage{mathblatt}
% im Vorspann definierter Baustein (fehlt in der Anleitung): Flächenbild eines Rechtecks
% mit geteilten Seiten #1+#2 (oben) und #3+#4 (links), Innenfelder #5..#8
\newcommand{\flaechenbild}[8]{%
\begin{tikzpicture}[x=1cm,y=1cm]
\draw[thick] (0,0) rectangle (4.5,-4.5);
\draw (3,0) -- (3,-4.5);
\draw (0,-3) -- (4.5,-3);
\node at (1.5,0.3) {$#1$};
\node at (3.75,0.3) {$#2$};
\node at (-0.4,-1.5) {$#3$};
\node at (-0.4,-3.75) {$#4$};
\node at (1.5,-1.5) {$#5$};
\node at (3.75,-1.5) {$#6$};
\node at (1.5,-3.75) {$#7$};
\node at (3.75,-3.75) {$#8$};
\end{tikzpicture}}

\newif\ifmitzone
\mitzonefalse
\begin{document}
\blattkopf{Terme und binomische Formeln}{Lernblatt}
\weit
\verzeichniszeile{\verz{e1}{1 Klammern auflösen (Nr. 10–16)} \verztrenn \verz{e2}{2 Ausklammern (Nr. 17–20)} \verztrenn \verz{e3}{3 Summe mal Summe (Nr. 21–28)} \verztrenn \verz{e4}{4 Binomische Formeln (Nr. 29–39)} \verztrenn \verz{e5}{5 Faktorisieren (Nr. 40–48)} \verztrenn \verz{abhaken}{Das kann ich}}

% Einheit 1 – Klammern auflösen (terme.md Einheit 3) – Nr. 10–16
\einheitenkopf[e1][1 Klammern auflösen]{Einheit 1 von 5 · Klammern auflösen}
\zweigzeile{Hier lernst du, Klammern in Termen aufzulösen · neu oder schon bekannt – je nach Buch · keine P10-Aufgabe · baut auf: Glieder zusammenfassen, Rechnen mit Vorzeichen}
\setcounter{aufgabe}{9}

\begin{aufgabe}{Ich erkenne, welches Vorzeichen zu welchem Glied gehört.}
\anweisung{Schreibe die Glieder mit ihrem Vorzeichen auf.}
\begin{geruest}
\gz{a}{$6x + 2y - 5$}{\leerfeld \quad \leerfeld \quad \leerfeld}
\gz{b}{$3x - 4 + y$}{\leerfeld \quad \leerfeld \quad \leerfeld}
\gz{c}{$-x + 7y - 2$}{\leerfeld \quad \leerfeld \quad \leerfeld}
\gz{d}{$8 - 3x - y$}{\leerfeld \quad \leerfeld \quad \leerfeld}
\gz{e}{$-2x - 5y + 9$}{\leerfeld \quad \leerfeld \quad \leerfeld}
\end{geruest}
\end{aufgabe}

\begin{aufgabe}{Ich erkenne, was vor einer Klammer steht.}
\anweisung{Kreuze an, was direkt vor der Klammer steht.}
\begin{teile}
\teil $5 + (x - 2)$ \hfill \kreuz{Plus} \kreuz{Minus} \kreuz{Zahl}
\teil $9x - (4 + x)$ \hfill \kreuz{Plus} \kreuz{Minus} \kreuz{Zahl}
\teil $3 \cdot (x + 1)$ \hfill \kreuz{Plus} \kreuz{Minus} \kreuz{Zahl}
\teil $2x - 4 \cdot (x - 3)$ \hfill \kreuz{Plus} \kreuz{Minus} \kreuz{Zahl}
\teil $7 - (2x - 1)$ \hfill \kreuz{Plus} \kreuz{Minus} \kreuz{Zahl}
\end{teile}
\end{aufgabe}

\begin{aufgabe}{Ich kann Klammern auflösen (neu in diesem Jahr).}
\anweisung{Löse die Klammer auf.}
\rechenplatz{2}
\begin{gleichungsraster}[1]
\gl{x + (y + 4)} & \gl{3x + (y - 2)} \\
\gl{5 + (2x - y)} & \gl{3 \cdot (x + 5)} \\
\gl{2 \cdot (y + 6)} & \gl{5 \cdot (x + 2)} \\
\gl{4 \cdot (y + 7)} & \gl{4x - (y + 2)} \\
\gl{x - (3y - z + 4)} & \gl{-3 \cdot (x - 4)} \\
\gl{x \cdot (x + 6)} & \\
\end{gleichungsraster}
\end{aufgabe}

\begin{aufgabe}{Ich kann Klammern auflösen und den Term zusammenfassen.}
\anweisung{Löse die Klammern auf und fasse zusammen.}
\rechenplatz{2}
\begin{gleichungsraster}[2]
\gl{3x + (2x + 4)} & \gl{6 + 2 \cdot (x + 1)} \\
\gl{5y + 3 \cdot (y - 2)} & \gl{7y - (2y + 5)} \\
\gl{10 - 2 \cdot (x - 4)} & \gl{2 \cdot (x + 3) + 4 \cdot (x - 1)} \\
\gl{3x \cdot (x - 2) - 2 \cdot (x^2 - 4x + 5)} & \\
\end{gleichungsraster}
\end{aufgabe}

\begin{aufgabe}{Ich finde den Fehler beim Auflösen einer Klammer.}
\anweisung{Finde den Fehler, benenne ihn und korrigiere die Rechnung.}
\begin{teile}
\teil Tom rechnet: \rechnung{9x - (4x - 6) &= 9x - 4x - 6 \\ &= 5x - 6}
\teil Mia rechnet: \rechnung{12 - 3 \cdot (x - 2) &= 12 - 3x - 6 \\ &= 6 - 3x}
\end{teile}
\end{aufgabe}

\begin{aufgabe}{Ich kann begründen, ob zwei Terme gleichwertig sind.}
\anweisung{Entscheide, ob die beiden Terme gleichwertig sind, und begründe.}
\begin{teile}
\teil $2 \cdot (x + 5)$ \ und \ $2x + 5$
\teil $10 - (x - 3)$ \ und \ $13 - x$
\teil $-(4 - x)$ \ und \ $x - 4$
\end{teile}
\end{aufgabe}

\begin{aufgabe}{Ich kann eine Sachaufgabe mit Klammern lösen.}
Ein rechteckiger Garten ist $(x + 6)$\,m lang und $x$\,m breit. Er bekommt rundherum einen Zaun, nur für das Tor bleiben $3$\,m frei.

\rechteck[seiten={x+6,x,,},punkte=]{4.5}{2.5}

\begin{teile}
\teil Stelle einen Term für die Länge des Zauns auf und vereinfache ihn.
\teil Der Garten ist $8$\,m breit. Wie viele Meter Zaun braucht man?
\end{teile}
\end{aufgabe}

\clearpage
% Einheit 2 – Ausklammern (terme.md Einheit 4) – Nr. 17–20
\einheitenkopf[e2][2 Ausklammern]{Einheit 2 von 5 · Ausklammern}
\zweigzeile{Hier lernst du, einen gemeinsamen Faktor auszuklammern · neu oder schon bekannt – je nach Buch · keine P10-Aufgabe · baut auf: Klammern auflösen (Einheit 1)}
\setcounter{aufgabe}{16}

\begin{aufgabe}{Ich kann einen gemeinsamen Faktor ausklammern.}
\anweisung{Klammere so viel wie möglich aus.}
\rechenplatz{2}
\begin{gleichungsraster}[1]
\gl{3x + 12} & \gl{5y + 35} \\
\gl{2x - 10} & \gl{7y + 21} \\
\gl{x^2 + 9x} & \gl{6x^2 + 15x} \\
\gl{6x + 6} & \gl{4x^2 + 8x + 12} \\
\anweisung{Klammere aus und mache die Probe durch Ausmultiplizieren.}
\gl[2]{10x^2 - 4x} & \gl[2]{12x^3 + 8x^2 - 4x} \\
\end{gleichungsraster}
\end{aufgabe}

\begin{aufgabe}{Ich finde den Fehler beim Ausklammern.}
\anweisung{Finde den Fehler, benenne ihn und korrigiere die Rechnung.}
\begin{teile}
\teil Jonas klammert aus: \rechnung{6x + 15 &= 3 \cdot (2x + 15)}
\teil Emma klammert aus: \rechnung{7x^2 + 7x &= 7x \cdot x}
\end{teile}
\end{aufgabe}

\begin{aufgabe}{Ich kann entscheiden, ob alle Glieder einen gemeinsamen Faktor haben.}
\anweisung{Haben alle Glieder einen gemeinsamen Faktor außer 1? Kreuze an, ohne zu rechnen.}
\begin{teilezwei}
\tz $4x + 9$ \quad \janein & \tz $6x + 10$ \quad \janein \\
\tz $x^2 + 3$ \quad \janein & \tz $5x^2 + 2x$ \quad \janein \\
\tz $3x + 7y$ \quad \janein & \tz $14y - 21$ \quad \janein \\
\end{teilezwei}
\end{aufgabe}

\begin{aufgabe}{Ich kann eine Sachaufgabe durch Ausklammern lösen.}
Ein rechteckiger Schulhof hat den Flächeninhalt $(8x^2 + 20x)$\,m$^2$. Er ist $4x$\,m lang.

\rechteck[seiten={4x,?,,},punkte=]{4.5}{2.5}

\begin{teile}
\teil Wie breit ist der Schulhof? Klammere dazu $4x$ aus.
\teil Für den Schulhof gilt $x = 10$. Wie lang und wie breit ist er? Prüfe mit dem Flächeninhalt.
\end{teile}
\end{aufgabe}

\clearpage
% Einheit 3 – Summe mal Summe (binomische-formeln.md Einheit 1) – Nr. 21–28
\einheitenkopf[e3][3 Summe mal Summe]{Einheit 3 von 5 · Summe mal Summe}
\zweigzeile{Hier lernst du, Terme mit zwei Variablen umzuformen und zwei Klammern miteinander malzunehmen · neu oder schon bekannt – je nach Buch · keine P10-Aufgabe · baut auf: Klammern auflösen (Einheit 1), Terme malnehmen}
\setcounter{aufgabe}{20}

\begin{aufgabe}{Ich erkenne, was mit was malgenommen wird.}
\anweisung{Zeichne Pfeile von jedem Glied der ersten Klammer zu jedem Glied der zweiten Klammer. Rechne nichts aus.}
\begin{teilezwei}
\tz \rule{0pt}{12mm}{\large $(x + 2) \cdot (y + 5)$} & \tz \rule{0pt}{12mm}{\large $(u + 4) \cdot (v - 1)$} \\
\tz \rule{0pt}{12mm}{\large $(x - 3) \cdot (x + 8)$} & \tz \rule{0pt}{12mm}{\large $(2x + 1) \cdot (x - 6)$} \\
\tz \rule{0pt}{12mm}{\large $(y + 7) \cdot (z + 3)$} & \\
\end{teilezwei}
\end{aufgabe}

\verfahren{Terme mit zwei Variablen}

\begin{aufgabe}{Ich kann den Wert eines Terms mit zwei Variablen berechnen.}
\anweisung{Setze ein und berechne den Wert des Terms.}
\begin{geruest}
\gz{a}{$2x + y$ \quad für $x = 3$ und $y = 4$}{\leerfeld}
\gz{b}{$x + 3y$ \quad für $x = 2$ und $y = 5$}{\leerfeld}
\gz{c}{$4x - y$ \quad für $x = 2$ und $y = 9$}{\leerfeld}
\gz{d}{$xy + 1$ \quad für $x = 3$ und $y = 4$}{\leerfeld}
\gz{e}{$3x - 2y$ \quad für $x = -2$ und $y = 5$}{\leerfeld}
\gz{f}{$x^2 + xy$ \quad für $x = -3$ und $y = 2$}{\leerfeld}
\gz{g}{$2x^2 - xy + y^2$ \quad für $x = -1$ und $y = 4$}{\leerfeld}
\end{geruest}
\end{aufgabe}

\begin{aufgabe}{Ich kann Terme mit zwei Variablen zusammenfassen (neu in diesem Jahr).}
\anweisung{Fasse zusammen.}
\begin{gleichungsraster}[1]
\gl{3x + 2y + 4x} & \gl{5u + v + 2v} \\
\gl{2x + 6y + 3x + 4y} & \gl{4u + 3v + 2u + 5v} \\
\gl{6x - y - 2x + 4y} & \gl{x^2 + 3xy + 2x^2 - xy + 5x} \\
\gl{4xy - 2x^2 + 3yx - x^2 + 6y} & \\
\end{gleichungsraster}
\end{aufgabe}

\begin{aufgabe}{Ich kann eine Klammer mit zwei Variablen auflösen.}
\anweisung{Löse die Klammer auf.}
\rechenplatz{2}
\begin{gleichungsraster}[1]
\gl{3 \cdot (x + 2y)} & \gl{4 \cdot (2x - y)} \\
\gl{5 \cdot (u + v)} & \gl{2 \cdot (3x + 4y)} \\
\gl{x \cdot (y + 3)} & \gl{-(2x - 5y)} \\
\gl{2x \cdot (x + 3y)} & \\
\anweisung{Löse die Klammern auf und fasse zusammen.}
\gl[2]{4 \cdot (x + 2y) - 3 \cdot (x - y)} & \\
\end{gleichungsraster}
\end{aufgabe}

\verfahren{Klammer mal Klammer}

\begin{aufgabe}{Ich kann zwei Klammern miteinander malnehmen.}
\anweisung{Multipliziere die Klammern aus.}
\rechenplatz{3}
\begin{gleichungsraster}[2]
\gl{(x + 2) \cdot (y + 4)} & \gl{(x + 5) \cdot (y + 1)} \\
\gl{(u + 3) \cdot (v + 6)} & \gl{(x + 7) \cdot (y + 2)} \\
\anweisung{Multipliziere aus und fasse zusammen.}
\gl{(x + 2) \cdot (x + 6)} & \gl{(x - 4) \cdot (x + 9)} \\
\gl{(x - 3) \cdot (x - 8)} & \gl{(2x + 1) \cdot (x + 7)} \\
\gl{(x + 2y) \cdot (3x + y)} & \gl{(3x - 2) \cdot (2x - 5)} \\
\end{gleichungsraster}
\end{aufgabe}

\begin{aufgabe}{Ich finde den Fehler beim Malnehmen zweier Klammern.}
\anweisung{Finde den Fehler, benenne ihn und korrigiere die Rechnung.}
\begin{teile}
\teil Paul rechnet: \rechnung{(x + 4) \cdot (x + 6) &= x^2 + 24}
\teil Lina rechnet: \rechnung{(x - 5) \cdot (x - 2) &= x^2 - 2x - 5x - 10 \\ &= x^2 - 7x - 10}
\teil Can rechnet: \rechnung{(x + 7) \cdot (x + 1) &= x^2 + x + 7x + 7 \\ &= x^2 + 15x}
\end{teile}
\end{aufgabe}

\begin{aufgabe}{Ich kann am Rechteck begründen, warum vier Produkte entstehen.}
Das Rechteck hat die Seitenlängen $x + 3$ und $x + 2$.

\flaechenbild{x}{3}{x}{2}{}{}{}{}

\begin{teile}
\teil Begründe am Bild, warum beim Malnehmen von $(x + 3) \cdot (x + 2)$ vier Produkte entstehen.
\teil Tim sagt: „$(x + 3) \cdot (x + 2) = x^2 + 6$.“ Begründe am Bild, was er vergessen hat.
\end{teile}
\end{aufgabe}

\begin{aufgabe}{Ich kann eine Sachaufgabe mit zwei Klammern lösen.}
Ein quadratisches Beet hat die Seitenlänge $x$\,m. Es wird neu angelegt: $3$\,m länger und $2$\,m schmaler.

\rechteck[seiten={x+3,x-2,,},punkte=]{4.5}{2.5}

\begin{teile}
\teil Stelle einen Term für den Flächeninhalt des neuen Beets auf und multipliziere aus.
\teil Um wie viel m$^2$ ändert sich der Flächeninhalt gegenüber dem alten Beet? Gib einen Term an.
\teil Wird ein Beet mit $x = 4$ dadurch größer oder kleiner? Und eines mit $x = 10$?
\end{teile}
\end{aufgabe}

\clearpage
% Einheit 4 – Binomische Formeln (binomische-formeln.md Einheit 2) – Nr. 29–39
\einheitenkopf[e4][4 Binomische Formeln]{Einheit 4 von 5 · Binomische Formeln}
\zweigzeile{Hier lernst du, Quadrate von Klammern mit den binomischen Formeln auszumultiplizieren · neu in diesem Jahr · P10 · baut auf: Klammer mal Klammer (Einheit 3), Quadratzahlen}
\setcounter{aufgabe}{28}

\begin{aufgabe}{Ich erkenne, ob eine Klammer mit sich selbst malgenommen wird.}
\anweisung{Wird hier eine Klammer mit sich selbst malgenommen? Kreuze an.}
\begin{teilezwei}
\tz $(x + 6)^2$ \quad \janein & \tz $(x + 6) \cdot (x - 6)$ \quad \janein \\
\tz $(y - 2) \cdot (y - 2)$ \quad \janein & \tz $(z + 5) \cdot (5 + z)$ \quad \janein \\
\tz $(2x + 1) \cdot (x + 2)$ \quad \janein & \\
\end{teilezwei}
\end{aufgabe}

\begin{aufgabe}{Ich erkenne, welche binomische Formel passt.}
\anweisung{Kreuze an, welche binomische Formel passt. Rechne nichts aus.}
\begin{teile}
\teil $(x + 8)^2$ \hfill \kreuz{erste} \kreuz{zweite} \kreuz{dritte} \kreuz{keine}
\teil $(y - 9)^2$ \hfill \kreuz{erste} \kreuz{zweite} \kreuz{dritte} \kreuz{keine}
\teil $(x + 4) \cdot (x - 4)$ \hfill \kreuz{erste} \kreuz{zweite} \kreuz{dritte} \kreuz{keine}
\teil $(x + 2) \cdot (x + 7)$ \hfill \kreuz{erste} \kreuz{zweite} \kreuz{dritte} \kreuz{keine}
\teil $(z - 1) \cdot (z - 1)$ \hfill \kreuz{erste} \kreuz{zweite} \kreuz{dritte} \kreuz{keine}
\end{teile}
\end{aufgabe}

\begin{aufgabe}{Ich erkenne, was $a$ und was $b$ ist.}
\anweisung{In den Formeln steht $a$ für das erste und $b$ für das zweite Glied in der Klammer. Schreibe $a$ und $b$ auf. Rechne nichts aus.}
\begin{geruest}
\gz{a}{$(x + 9)^2$}{\feld{a} \quad \feld{b}}
\gz{b}{$(y - 4)^2$}{\feld{a} \quad \feld{b}}
\gz{c}{$(3x + 1)^2$}{\feld{a} \quad \feld{b}}
\gz{d}{$(5 + 2y)^2$}{\feld{a} \quad \feld{b}}
\gz{e}{$(4x - 7y)^2$}{\feld{a} \quad \feld{b}}
\end{geruest}
\end{aufgabe}

\verfahren{Ausmultiplizieren mit den Formeln}

\begin{aufgabe}{Ich kann mit den binomischen Formeln ausmultiplizieren.}
\anweisung{Multipliziere aus.}
\rechenplatz{2}
\begin{gleichungsraster}[2]
\gl{(x + 4)^2} & \gl{(x + 1)^2} \\
\gl{(y + 6)^2} & \gl{(x + 10)^2} \\
\gl{(x - 9)^2} & \gl{(x + 11) \cdot (x - 11)} \\
\gl{(x + 3) \cdot (x - 8)} & \\
\end{gleichungsraster}
\end{aufgabe}

\begin{aufgabe}{Ich kann mit den binomischen Formeln ausmultiplizieren – weiter.}
\begin{gleichungsraster}[2]
\gl{(3x + 4)^2} & \gl{(x - 5y)^2} \\
\gl{2 \cdot (x + 6)^2} & \gl{-(x - 8)^2} \\
\gl{(y - 4)^2 - 10} & \gl{(x - 6)^2 - 2 \cdot (3x - 4) + 7x \quad \text{(P10 2022 GYM)}} \\
\end{gleichungsraster}
\end{aufgabe}

\verfahren{Zwei Terme als gleich nachweisen}

\begin{aufgabe}{Ich kann nachweisen, dass zwei Terme gleich sind.}
\anweisung{Zeige durch Umformen, dass beide Terme gleich sind.}
\rechenplatz{2}
\begin{gleichungsraster}[2]
\gl{(x + 2)^2 - 1 \ \text{ und } \ x^2 + 4x + 3} & \gl{(x + 1)^2 + 7 \ \text{ und } \ x^2 + 2x + 8} \\
\gl{(x - 4)^2 + 2 \ \text{ und } \ x^2 - 8x + 18} & \\
\anweisung{Eine Parabel hat die Gleichung $y = (x + 6)^2 - 11$. Zeige, dass sie auch die Gleichung $y = x^2 + 12x + 25$ hat. (P10 2017 OS)}
\gl[3]{y = (x + 6)^2 - 11} & \\
\end{gleichungsraster}
\end{aufgabe}

\verfahren{Kopfrechnen mit den Formeln}

\begin{aufgabe}{Ich kann mit den binomischen Formeln im Kopf rechnen.}
\anweisung{Rechne mit einer binomischen Formel im Kopf.}
\begin{teilezwei}
\tz $21^2 = $ \leerfeld & \tz $51^2 = $ \leerfeld \\
\tz $19^2 = $ \leerfeld & \tz $98^2 = $ \leerfeld \\
\tz $32 \cdot 28 = $ \leerfeld & \tz $45 \cdot 55 = $ \leerfeld \\
\end{teilezwei}
\end{aufgabe}

\begin{aufgabe}{Ich finde den Fehler beim Ausmultiplizieren mit einer binomischen Formel.}
\anweisung{Finde den Fehler, benenne ihn und korrigiere die Rechnung.}
\begin{teile}
\teil Noah rechnet: \rechnung{(x + 6)^2 &= x^2 + 36}
\teil Ella rechnet: \rechnung{(3x - 2)^2 &= 3x^2 - 12x + 4}
\teil Finn rechnet: \rechnung{-(x + 9)^2 &= -x^2 + 18x + 81}
\end{teile}
\end{aufgabe}

\begin{aufgabe}{Ich kann entscheiden, welcher Term gleichwertig ist.}
\anweisung{Kreuze den Term an, der gleichwertig ist.}
\begin{teile}
\teil $(y + 9)^2$ \\
\kreuz{$y^2 + 81$} \\ \kreuz{$y^2 + 18y + 81$} \\ \kreuz{$y^2 + 9y + 81$} \\ \kreuz{$y^2 + 18y + 18$}
\teil $(x - 10) \cdot (x + 10)$ \\
\kreuz{$x^2 - 20x - 100$} \\ \kreuz{$x^2 - 100$} \\ \kreuz{$x^2 + 100$} \\ \kreuz{$x^2 - 20x + 100$}
\teil $(4 - 3y)^2$ \quad (P10 2019 GYM) \\
\kreuz{$16 - 9y^2$} \\ \kreuz{$16 - 24y + 9y^2$} \\ \kreuz{$16 - 12y + 9y^2$} \\ \kreuz{$16 - 24y - 9y^2$}
\end{teile}
\end{aufgabe}

\begin{aufgabe}{Ich kann begründen, warum beim Quadrat einer Summe ein Mittelglied entsteht.}
\begin{teile}
\teil Rechne $(6 + 2)^2$ und $6^2 + 2^2$. Begründe damit, warum man beim Quadrat einer Summe nicht jedes Glied einzeln quadrieren darf.
\teil Das Quadrat hat die Seitenlänge $x + 2$. Begründe am Bild, warum $(x + 2)^2$ nicht $x^2 + 4$ ist.
\end{teile}

\flaechenbild{x}{2}{x}{2}{}{}{}{}

\end{aufgabe}

\begin{aufgabe}{Ich kann eine Sachaufgabe mit einer binomischen Formel lösen.}
Ein quadratischer Spielplatz hat die Seitenlänge $x$\,m. Er wird nach rechts und nach unten um je $12$\,m vergrößert, sodass wieder ein Quadrat entsteht.

\flaechenbild{x}{12}{x}{12}{}{}{}{}

\begin{teile}
\teil Stelle einen Term für den Flächeninhalt des neuen Spielplatzes auf und multipliziere aus.
\teil Um wie viel m$^2$ wird der Spielplatz größer? Gib einen Term an.
\teil Der alte Spielplatz war $20$\,m lang. Wie viel m$^2$ kommen hinzu?
\end{teile}
\end{aufgabe}

\clearpage
% Einheit 5 – Faktorisieren (binomische-formeln.md Einheit 3) – Nr. 40–48
\einheitenkopf[e5][5 Faktorisieren]{Einheit 5 von 5 · Faktorisieren}
\zweigzeile{Hier lernst du, eine Summe mit den binomischen Formeln in ein Produkt zu verwandeln · neu in diesem Jahr · keine P10-Aufgabe · baut auf: Ausklammern (Einheit 2), binomische Formeln (Einheit 4), Quadratwurzeln}
\setcounter{aufgabe}{39}

\begin{aufgabe}{Ich erkenne, ob ein Term ein Quadrat ist.}
\anweisung{Ist der Term ein Quadrat? Kreuze an.}
\begin{teilezwei}
\tz $y^2$ \quad \janein & \tz $64$ \quad \janein \\
\tz $25x^2$ \quad \janein & \tz $12x$ \quad \janein \\
\tz $18$ \quad \janein & \tz $3x^2$ \quad \janein \\
\tz $100y^2$ \quad \janein & \\
\end{teilezwei}
\end{aufgabe}

\begin{aufgabe}{Ich erkenne die beiden Quadrate und das Mittelglied.}
\anweisung{Unterstreiche die beiden Quadrate und kreise das Mittelglied ein. Rechne nichts aus.}
\begin{teilezwei}
\tz \rule{0pt}{9mm}{\large $y^2 + 16y + 64$} & \tz \rule{0pt}{9mm}{\large $36 + 12x + x^2$} \\
\tz \rule{0pt}{9mm}{\large $9x^2 - 24x + 16$} & \tz \rule{0pt}{9mm}{\large $49 - 14y + y^2$} \\
\tz \rule{0pt}{9mm}{\large $81x^2 + 18x + 1$} & \\
\end{teilezwei}
\end{aufgabe}

\verfahren{Faktorisieren mit den Formeln}

\begin{aufgabe}{Ich kann einen Term mit einer binomischen Formel faktorisieren.}
\anweisung{Faktorisiere.}
\rechenplatz{2}
\begin{gleichungsraster}[1]
\gl{x^2 - 16} & \gl{x^2 - 36} \\
\gl{y^2 - 100} & \gl{x^2 - 1} \\
\gl{x^2 + 12x + 36} & \gl{x^2 - 16x + 64} \\
\gl{9x^2 + 30x + 25} & \gl[2]{2x^2 - 72} \\
\anweisung{Faktorisiere und mache die Probe durch Ausmultiplizieren.}
\gl[2]{x^2 - 20x + 100} & \gl[3]{5x^3 - 40x^2 + 80x} \\
\end{gleichungsraster}
\end{aufgabe}

\verfahren{Gleichungen durch Faktorisieren lösen}

\begin{aufgabe}{Ich kann eine Gleichung lösen, indem ich faktorisiere.}
\anweisung{Faktorisiere und löse die Gleichung.}
\rechenplatz{3}
\begin{gleichungsraster}[3]
\gl{x^2 - 81 = 0} & \gl{x^2 - 144 = 0} \\
\gl{x^2 - 18x + 81 = 0} & \gl{3x^2 - 48 = 0} \\
\end{gleichungsraster}
\end{aufgabe}

\verfahren{Quadratische Ergänzung}

\begin{aufgabe}{Ich kann einen Term als Quadrat einer Klammer plus eine Zahl schreiben.}
\anweisung{Schreibe den Term als Quadrat einer Klammer plus oder minus eine Zahl.}
\rechenplatz{3}
\begin{gleichungsraster}[2]
\gl{x^2 + 4x + 7} & \gl{x^2 + 8x + 20} \\
\gl{x^2 - 12x + 40} & \gl{x^2 + 16x + 50} \\
\end{gleichungsraster}
\end{aufgabe}

\begin{aufgabe}{Ich finde den Fehler beim Faktorisieren.}
\anweisung{Finde den Fehler, benenne ihn und korrigiere.}
\begin{teile}
\teil Jana faktorisiert: \rechnung{x^2 + 64 &= (x + 8) \cdot (x - 8)}
\teil Ali faktorisiert: \rechnung{x^2 + 9x + 16 &= (x + 4)^2}
\teil Tim faktorisiert: \rechnung{2x^2 - 32 &= 2 \cdot (x^2 - 32)}
\end{teile}
\end{aufgabe}

\begin{aufgabe}{Ich kann entscheiden, ob sich ein Term mit einer binomischen Formel faktorisieren lässt.}
\anweisung{Lässt sich der Term mit einer binomischen Formel faktorisieren? Kreuze an, ohne zu rechnen.}
\begin{teilezwei}
\tz $x^2 + 22x + 121$ \quad \janein & \tz $x^2 + 36$ \quad \janein \\
\tz $y^2 - 169$ \quad \janein & \tz $x^2 + 5x + 4$ \quad \janein \\
\tz $4x^2 - 20x + 25$ \quad \janein & \tz $x^2 - 8x - 16$ \quad \janein \\
\end{teilezwei}
\end{aufgabe}

\begin{aufgabe}{Ich kann begründen, wann ein Term keine binomische Formel ist.}
\begin{teile}
\teil Begründe, warum sich $x^2 - 100$ mit einer binomischen Formel faktorisieren lässt, $x^2 + 100$ aber nicht.
\teil $x^2$ und $36$ sind Quadrate. Begründe, warum $x^2 + 10x + 36$ trotzdem keine binomische Formel ist.
\end{teile}
\end{aufgabe}

\begin{aufgabe}{Ich kann eine Sachaufgabe durch Faktorisieren lösen.}
Eine quadratische Terrasse hat den Flächeninhalt $(x^2 + 18x + 81)$\,m$^2$.

\rechteck[seiten={?,?,,},punkte=]{3}{3}

\begin{teile}
\teil Wie lang ist eine Seite der Terrasse? Faktorisiere dazu.
\teil Gib einen Term für den Umfang der Terrasse an.
\teil Für die Terrasse gilt $x = 3$. Wie lang ist eine Seite, und wie groß ist die Terrasse?
\end{teile}
\end{aufgabe}

\begin{abhakseite}
\ifmitzone
\abhakgruppe{Kennst du schon}
\abhak{1}{Ich kann negative Zahlen addieren und subtrahieren.}
\abhak{2}{Ich kann Zahlen mit Vorzeichen malnehmen.}
\abhak{3}{Ich kann Punkt vor Strich rechnen.}
\abhak{4}{Ich kann den Wert eines Terms berechnen.}
\abhak{5}{Ich kann gleichartige Glieder zusammenfassen.}
\abhak{6}{Ich kann Terme malnehmen.}
\abhak{7}{Ich finde den Fehler beim Zusammenfassen.}
\abhak{8}{Ich kann Terme mit $x$ und $x^2$ richtig zusammenfassen.}
\abhak{9}{Ich kann Quadratzahlen und Quadratwurzeln angeben.}
\fi
\abhakgruppe{Einheit 1 · Klammern auflösen}
\abhak{10}{Ich erkenne, welches Vorzeichen zu welchem Glied gehört.}
\abhak{11}{Ich erkenne, was vor einer Klammer steht.}
\abhak{12}{Ich kann Klammern auflösen (neu in diesem Jahr).}
\abhak{13}{Ich kann Klammern auflösen und den Term zusammenfassen.}
\abhak{14}{Ich finde den Fehler beim Auflösen einer Klammer.}
\abhak{15}{Ich kann begründen, ob zwei Terme gleichwertig sind.}
\abhak{16}{Ich kann eine Sachaufgabe mit Klammern lösen.}
\abhakgruppe{Einheit 2 · Ausklammern}
\abhak{17}{Ich kann einen gemeinsamen Faktor ausklammern.}
\abhak{18}{Ich finde den Fehler beim Ausklammern.}
\abhak{19}{Ich kann entscheiden, ob alle Glieder einen gemeinsamen Faktor haben.}
\abhak{20}{Ich kann eine Sachaufgabe durch Ausklammern lösen.}
\abhakgruppe{Einheit 3 · Summe mal Summe}
\abhak{21}{Ich erkenne, was mit was malgenommen wird.}
\abhakgruppe{\normalfont\itshape Terme mit zwei Variablen}
\abhak{22}{Ich kann den Wert eines Terms mit zwei Variablen berechnen.}
\abhak{23}{Ich kann Terme mit zwei Variablen zusammenfassen (neu in diesem Jahr).}
\abhak{24}{Ich kann eine Klammer mit zwei Variablen auflösen.}
\abhakgruppe{\normalfont\itshape Klammer mal Klammer}
\abhak{25}{Ich kann zwei Klammern miteinander malnehmen.}
\abhak{26}{Ich finde den Fehler beim Malnehmen zweier Klammern.}
\abhak{27}{Ich kann am Rechteck begründen, warum vier Produkte entstehen.}
\abhak{28}{Ich kann eine Sachaufgabe mit zwei Klammern lösen.}
\abhakgruppe{Einheit 4 · Binomische Formeln}
\abhak{29}{Ich erkenne, ob eine Klammer mit sich selbst malgenommen wird.}
\abhak{30}{Ich erkenne, welche binomische Formel passt.}
\abhak{31}{Ich erkenne, was $a$ und was $b$ ist.}
\abhakgruppe{\normalfont\itshape Ausmultiplizieren mit den Formeln}
\abhak{32}{Ich kann mit den binomischen Formeln ausmultiplizieren.}
\abhak{33}{Ich kann mit den binomischen Formeln ausmultiplizieren – weiter.}
\abhakgruppe{\normalfont\itshape Zwei Terme als gleich nachweisen}
\abhak{34}{Ich kann nachweisen, dass zwei Terme gleich sind.}
\abhakgruppe{\normalfont\itshape Kopfrechnen mit den Formeln}
\abhak{35}{Ich kann mit den binomischen Formeln im Kopf rechnen.}
\abhak{36}{Ich finde den Fehler beim Ausmultiplizieren mit einer binomischen Formel.}
\abhak{37}{Ich kann entscheiden, welcher Term gleichwertig ist.}
\abhak{38}{Ich kann begründen, warum beim Quadrat einer Summe ein Mittelglied entsteht.}
\abhak{39}{Ich kann eine Sachaufgabe mit einer binomischen Formel lösen.}
\abhakgruppe{Einheit 5 · Faktorisieren}
\abhak{40}{Ich erkenne, ob ein Term ein Quadrat ist.}
\abhak{41}{Ich erkenne die beiden Quadrate und das Mittelglied.}
\abhakgruppe{\normalfont\itshape Faktorisieren mit den Formeln}
\abhak{42}{Ich kann einen Term mit einer binomischen Formel faktorisieren.}
\abhakgruppe{\normalfont\itshape Gleichungen durch Faktorisieren lösen}
\abhak{43}{Ich kann eine Gleichung lösen, indem ich faktorisiere.}
\abhakgruppe{\normalfont\itshape Quadratische Ergänzung}
\abhak{44}{Ich kann einen Term als Quadrat einer Klammer plus eine Zahl schreiben.}
\abhak{45}{Ich finde den Fehler beim Faktorisieren.}
\abhak{46}{Ich kann entscheiden, ob sich ein Term mit einer binomischen Formel faktorisieren lässt.}
\abhak{47}{Ich kann begründen, wann ein Term keine binomische Formel ist.}
\abhak{48}{Ich kann eine Sachaufgabe durch Faktorisieren lösen.}
\end{abhakseite}

\end{document}
